% Options for packages loaded elsewhere
\PassOptionsToPackage{unicode}{hyperref}
\PassOptionsToPackage{hyphens}{url}
\documentclass[
  spanish,
  11pt,
]{article}
\usepackage{xcolor}
\usepackage[margin=2.5cm]{geometry}
\usepackage{amsmath,amssymb}
\setcounter{secnumdepth}{-\maxdimen} % remove section numbering
\usepackage{iftex}
\ifPDFTeX
  \usepackage[T1]{fontenc}
  \usepackage[utf8]{inputenc}
  \usepackage{textcomp} % provide euro and other symbols
\else % if luatex or xetex
  \usepackage{unicode-math} % this also loads fontspec
  \defaultfontfeatures{Scale=MatchLowercase}
  \defaultfontfeatures[\rmfamily]{Ligatures=TeX,Scale=1}
\fi
\usepackage{lmodern}
\ifPDFTeX\else
  % xetex/luatex font selection
  \setmainfont[]{Arial}
\fi
% Use upquote if available, for straight quotes in verbatim environments
\IfFileExists{upquote.sty}{\usepackage{upquote}}{}
\IfFileExists{microtype.sty}{% use microtype if available
  \usepackage[]{microtype}
  \UseMicrotypeSet[protrusion]{basicmath} % disable protrusion for tt fonts
}{}
\usepackage{setspace}
\makeatletter
\@ifundefined{KOMAClassName}{% if non-KOMA class
  \IfFileExists{parskip.sty}{%
    \usepackage{parskip}
  }{% else
    \setlength{\parindent}{0pt}
    \setlength{\parskip}{6pt plus 2pt minus 1pt}}
}{% if KOMA class
  \KOMAoptions{parskip=half}}
\makeatother
\usepackage{longtable,booktabs,array}
\newcounter{none} % for unnumbered tables
\usepackage{calc} % for calculating minipage widths
% Correct order of tables after \paragraph or \subparagraph
\usepackage{etoolbox}
\makeatletter
\patchcmd\longtable{\par}{\if@noskipsec\mbox{}\fi\par}{}{}
\makeatother
% Allow footnotes in longtable head/foot
\IfFileExists{footnotehyper.sty}{\usepackage{footnotehyper}}{\usepackage{footnote}}
\makesavenoteenv{longtable}
\usepackage{graphicx}
\makeatletter
\newsavebox\pandoc@box
\newcommand*\pandocbounded[1]{% scales image to fit in text height/width
  \sbox\pandoc@box{#1}%
  \Gscale@div\@tempa{\textheight}{\dimexpr\ht\pandoc@box+\dp\pandoc@box\relax}%
  \Gscale@div\@tempb{\linewidth}{\wd\pandoc@box}%
  \ifdim\@tempb\p@<\@tempa\p@\let\@tempa\@tempb\fi% select the smaller of both
  \ifdim\@tempa\p@<\p@\scalebox{\@tempa}{\usebox\pandoc@box}%
  \else\usebox{\pandoc@box}%
  \fi%
}
% Set default figure placement to htbp
\def\fps@figure{htbp}
\makeatother
\ifLuaTeX
\usepackage[bidi=basic,shorthands=off]{babel}
\else
\usepackage[bidi=default,shorthands=off]{babel}
\fi
\ifPDFTeX
\else
\babelfont{rm}[]{Arial}
\fi
\ifLuaTeX
  \usepackage{selnolig} % disable illegal ligatures
\fi
\setlength{\emergencystretch}{3em} % prevent overfull lines
\providecommand{\tightlist}{%
  \setlength{\itemsep}{0pt}\setlength{\parskip}{0pt}}
\usepackage{bookmark}
\IfFileExists{xurl.sty}{\usepackage{xurl}}{} % add URL line breaks if available
\urlstyle{same}
% fallback for those not using the hyperref driver hyperxmp:
\makeatletter
\@ifundefined{xmpquote}{\newcommand{\xmpquote}[1]{#1}}{}
\makeatother
\hypersetup{
  pdftitle={Pronóstico mensual de las órdenes entrantes de LightForce Orthodontics mediante series temporales},
  pdflang={es},
  hidelinks,
  pdfcreator={LaTeX via pandoc}}

\title{Pronóstico mensual de las órdenes entrantes de LightForce
Orthodontics mediante series temporales}
\author{}
\date{}

\begin{document}
\maketitle

\setstretch{1.5}
\textbf{{[}Insertar logo de Lead University{]}}

\textbf{Curso:} Modelado matemático (BCD5105), III Cuatrimestre 2026

\textbf{Proyecto integrador:} Avance 1, planteamiento y primer modelo

\textbf{Título:} Pronóstico mensual de las órdenes entrantes de
LightForce Orthodontics mediante series temporales

\textbf{Grupo:} \emph{(número)}

\textbf{Integrantes:} \emph{(nombres completos)}

\textbf{Fecha de entrega:} 29 de octubre de 2026

\textbf{Notebook de Colab (permiso de lectura):} \emph{(pegar el enlace
compartido)}

\newpage

\section{1. Introducción y planteamiento del
problema}\label{introducciuxf3n-y-planteamiento-del-problema}

LightForce Orthodontics es una empresa de EE. UU. (fundada en 2019) que
fabrica aparatos ortodónticos totalmente personalizados con impresión 3D
y CAD/CAM, y que vende solo a ortodoncistas. Cada pedido nace cuando el
doctor envía el escaneo digital del paciente por el portal; el caso se
registra en el sistema de manufactura y pasa por segmentación,
planificación del tratamiento, aprobación del doctor, diseño, impresión
y envío. Como cada caso es único, el trabajo no se puede inventariar: la
carga de manufactura depende de cuántos escaneos llegan.

El problema que se aborda es \textbf{pronosticar cuántas órdenes
entrantes (escaneos recibidos) llegarán cada mes durante los próximos
tres meses}. Una orden entrante es un caso nuevo con escaneo recibido;
los reórdenes quedan fuera. El volumen mensual ronda las 4,700 órdenes
desde 2024, pero en 2026 se movió entre 4,095 y 6,032 por el lanzamiento
de un producto nuevo, el bracket de acero inoxidable: se habilitó de
forma progresiva, la demanda superó la capacidad de las máquinas y se
detuvo su recepción para ajustar la capacidad durante el último
trimestre (información de operaciones; los datos de la sección 4 son
consistentes con ella).

El contexto costarricense es directo. Tres de los cuatro supervisores de
manufactura digital trabajan en Costa Rica y desde aquí se atiende a los
clientes de EE. UU. y Canadá. La demanda la determinan los doctores del
norte, pero la capacidad la aporta un equipo en Costa Rica, de modo que
anticipar el volumen mensual es una entrada directa para planificar
turnos, horas extra y vacaciones. A mediano plazo el objetivo es un
pronóstico diario; este avance parte de la escala mensual, más simple y
apropiada para construir un primer modelo con los métodos del curso.

\section{2. Justificación}\label{justificaciuxf3n}

Un pronóstico de tres meses ayuda a decidir con tiempo cuánta capacidad
se necesita. Si la demanda supera lo previsto, los casos esperan en cola
y se alargan los tiempos de ciclo (SLA) que la empresa monitorea; si
queda por debajo, hay horas pagadas sin trabajo. Además, las cuentas de
cada doctor están asignadas a un equipo específico (Prism), por lo que
un exceso de casos no se reparte automáticamente entre equipos y
conviene anticiparlo.

Se benefician los supervisores y planificadores de manufactura digital,
que toman las decisiones de dotación; los técnicos en Costa Rica, que
enfrentan menos picos y valles; y los doctores, que reciben antes su
plan de tratamiento. El impacto económico no se cuantifica en este
avance; lo que se mide es cuánto mejora el pronóstico respecto de una
regla simple y cuándo falla.

\section{3. Objetivos}\label{objetivos}

\textbf{Objetivo general.} Modelar y pronosticar el número mensual de
órdenes entrantes de LightForce Orthodontics entre enero de 2023 y
septiembre de 2026 con un modelo SARIMA, y evaluar si supera de forma
consistente al pronóstico ingenuo estacional, para apoyar la
planificación de capacidad con un horizonte de tres meses.

\textbf{Objetivos específicos.}

\begin{enumerate}
\def\labelenumi{\arabic{enumi}.}
\tightlist
\item
  Construir la serie mensual de órdenes entrantes (45 meses) a partir de
  la tabla de órdenes de manufactura y documentar su definición.
\item
  Caracterizar la tendencia, la estacionalidad anual y el efecto del
  lanzamiento del bracket de acero inoxidable mediante estadísticas
  descriptivas, autocorrelación y descomposición STL.
\item
  Ajustar un SARIMA con período 12 y compararlo con el ingenuo
  estacional y AutoETS mediante MAE, RMSE y MASE, con validación de
  ventana expansiva (19 ventanas de tres meses).
\item
  Diagnosticar los residuos e identificar las limitaciones del modelo,
  para definir el segundo método de la entrega final.
\end{enumerate}

\section{4. Dataset y análisis
exploratorio}\label{dataset-y-anuxe1lisis-exploratorio}

\textbf{Fuente.} Base de datos analítica de LightForce Orthodontics
(Databricks), tabla \texttt{analytics.manufacturing.mfg\_all\_orders},
consultada el 26 de octubre de 2026 con autorización de la empresa. Es
una fuente interna, sin enlace público; la consulta SQL y la serie
resultante (45 filas) están en el notebook. Cada fila de la tabla es un
caso de manufactura.

\textbf{Definición y período.} Una orden entrante es un caso nuevo
(\texttt{order\_type\ =\ \textquotesingle{}New\ case\textquotesingle{}})
con fecha de recepción del escaneo (\texttt{scan\_received\_date}),
pasada de UTC a hora del este de EE. UU. antes de asignar el mes. Se
cuentan todos los estados, incluso casos luego cancelados, porque la
demanda llegó. El período es enero de 2023 a septiembre de 2026 (45
meses, 207,214 órdenes). Se excluyen 2019 a 2022, años de rampa de
crecimiento (66 casos en 2019 y 29,139 en 2022; desde 2023 el volumen
anual está entre 49,000 y 57,000). Octubre de 2026 solo tiene datos
hasta el día 25 (3,036 órdenes), por lo que no se usa en el ajuste.

{\def\LTcaptype{none} % do not increment counter
\begin{longtable}[]{@{}
  >{\raggedright\arraybackslash}p{(\linewidth - 6\tabcolsep) * \real{0.2500}}
  >{\raggedright\arraybackslash}p{(\linewidth - 6\tabcolsep) * \real{0.2500}}
  >{\raggedright\arraybackslash}p{(\linewidth - 6\tabcolsep) * \real{0.2500}}
  >{\raggedright\arraybackslash}p{(\linewidth - 6\tabcolsep) * \real{0.2500}}@{}}
\toprule\noalign{}
\begin{minipage}[b]{\linewidth}\raggedright
Variable
\end{minipage} & \begin{minipage}[b]{\linewidth}\raggedright
Tipo
\end{minipage} & \begin{minipage}[b]{\linewidth}\raggedright
Unidad
\end{minipage} & \begin{minipage}[b]{\linewidth}\raggedright
Descripción
\end{minipage} \\
\midrule\noalign{}
\endhead
\bottomrule\noalign{}
\endlastfoot
\texttt{mes} & Fecha (índice) & Mes calendario, hora ET & Mes en que se
recibió el escaneo \\
\texttt{ordenes} & Entero, no negativo & Órdenes por mes & Casos nuevos
con escaneo recibido (variable a pronosticar) \\
\texttt{acero} & Entero, no negativo & Órdenes por mes & De ellas, las
de bracket de acero inoxidable (\texttt{bracket\_material\_name}) \\
\texttt{dias\_habiles} & Entero (derivada) & Días, lunes a viernes &
Días hábiles del mes, entre 20 y 23 \\
\end{longtable}
}

Tabla 1. Variables del conjunto de datos.

\textbf{Estadísticas descriptivas.} La media mensual es 4,605 órdenes
(mediana 4,734; desviación 547; coeficiente de variación 0.12), con
mínimo de 3,294 (febrero de 2023) y máximo de 6,032 (junio de 2026). Los
totales anuales son 49,084 (2023; 4,090 por mes), 55,879 (2024; 4,657),
56,648 (2025; 4,721) y 45,603 de enero a septiembre de 2026 (5,067 por
mes).

\begin{figure}
\centering
\includegraphics[width=0.82\linewidth,height=\textheight,keepaspectratio,alt={Figura 1. Órdenes mensuales según el tipo de bracket.}]{./media-564a31000ad1b7c6/figuras/fig1_serie_mensual.png}
\caption{Figura 1. Órdenes mensuales según el tipo de bracket.}
\end{figure}

\textbf{Figura 1.} Tras crecer durante 2023, el volumen se estabiliza
desde 2024 entre 3,900 y 5,200 órdenes por mes. En 2026 aparece el
bracket de acero: 63 órdenes en febrero, 352 en abril, 1,625 en junio
(27 \% del mes) y 1,829 en julio (32 \%); luego cae a 117 en septiembre
(3 \%). El total de junio a agosto promedia 5,735, un 22 \% sobre el
promedio de octubre de 2025 a abril de 2026 (4,683), y septiembre
(4,095) queda 13 \% debajo. Las órdenes del resto de brackets oscilan
entre 3,721 y 4,407 de mayo a septiembre de 2026, bajo el rango de 4,145
a 5,158 de los diez meses previos, lo que sugiere que parte del acero
sustituyó a otros productos (hipótesis que no se prueba aquí).
\textbf{Implicación:} el auge y la caída de 2026 son una decisión de
capacidad, no una variación natural de la demanda, y los modelos deben
evaluarse por separado en ese tramo. La serie no se recorta ni se
ajusta: son órdenes realmente recibidas.

\begin{figure}
\centering
\includegraphics[width=0.82\linewidth,height=\textheight,keepaspectratio,alt={Figura 2. Perfil anual: un trazo por año.}]{./media-564a31000ad1b7c6/figuras/fig2_perfil_anual.png}
\caption{Figura 2. Perfil anual: un trazo por año.}
\end{figure}

\textbf{Figura 2.} Los trazos de 2024 y 2025 son casi iguales: los meses
más bajos son noviembre (índice 0.89 respecto del promedio de su año) y
diciembre (0.86), y los más altos enero, junio, julio y octubre (1.05 a
1.08). 2023 queda más abajo por el crecimiento, y 2026 se separa a
partir de mayo. \textbf{Implicación:} hay un patrón anual repetible pero
de amplitud moderada (cerca de ±10 \%), por lo que repetir el mismo mes
del año anterior será una referencia difícil de superar.

\begin{figure}
\centering
\includegraphics[width=0.82\linewidth,height=\textheight,keepaspectratio,alt={Figura 3. Descomposición STL (período 12, componente anual fijo).}]{./media-564a31000ad1b7c6/figuras/fig3_stl.png}
\caption{Figura 3. Descomposición STL (período 12, componente anual
fijo).}
\end{figure}

\textbf{Figura 3.} La fuerza de la tendencia es 0.65 y la de la
estacionalidad anual 0.51 (escala de 0 a 1), es decir, moderadas. La
tendencia sube en 2023, se aplana en torno a 4,700 y vuelve a subir en
2026. El resto más grande es septiembre de 2026 (−1,021 órdenes), la
caída del acero. La autocorrelación es 0.50 en el rezago 1, 0.35 en el 2
y solo 0.16 en el 12. \textbf{Implicación:} la dependencia entre meses
vecinos es más clara que la estacional, y con solo 45 meses (menos de
cuatro ciclos) la estacionalidad se estima con poca precisión. Las
pruebas de raíz unitaria no coinciden sobre la serie original (ADF p =
0.002; KPSS p = 0.02), pero sí tras la diferencia estacional (ADF p =
0.006; KPSS p ≥ 0.10).

\section{5. Formulación matemática
inicial}\label{formulaciuxf3n-matemuxe1tica-inicial}

\textbf{Variables.} Sea \(y_t \in \mathbb{N}\) el número de órdenes
entrantes del mes \(t\), con \(t = 1, \dots, T\) y \(T = 45\). Se busca
el pronóstico \(\hat y_{T+h \mid T}\) para \(h = 1, 2, 3\). El período
estacional es \(m = 12\) y \(B\) es el operador de rezago
(\(B y_t = y_{t-1}\)).

\textbf{Modelo.} Un SARIMA\((p,d,q)(P,D,Q)_m\) cumple

\[\phi(B)\,\Phi(B^m)\,(1-B)^d(1-B^m)^D\, y_t = c + \theta(B)\,\Theta(B^m)\,\varepsilon_t, \qquad \varepsilon_t \sim (0, \sigma^2),\]

donde \(\phi\), \(\theta\), \(\Phi\) y \(\Theta\) son polinomios de
grados \(p\), \(q\), \(P\) y \(Q\). Los órdenes se identifican con
AutoARIMA (pruebas de raíz unitaria para \(d\) y \(D\), y criterio AICc
para el resto) y los coeficientes se estiman por máxima verosimilitud.
Con los 45 meses se obtiene SARIMA\((1,1,0)(0,1,0)_{12}\), es decir,

\[(1-\phi_1 B)(1-B)(1-B^{12})\,y_t = \varepsilon_t.\]

En palabras: el cambio de un mes al siguiente, medido contra el mismo
cambio un año antes, depende del cambio del mes anterior.

\textbf{Restricciones.} (i) \(y_t \ge 0\): todo pronóstico se trunca en
cero, \(\hat y^{+}_{t} = \max(0, \hat y_t)\). (ii) Estacionariedad e
invertibilidad: las raíces de los polinomios quedan fuera del círculo
unitario. (iii) Validación sin fuga: cada pronóstico usa solo datos
anteriores a su origen, y el orden se vuelve a identificar en cada
ventana.

\textbf{Supuestos y su justificación.}

\begin{itemize}
\tightlist
\item
  \emph{Período anual fijo (\(m = 12\)).} Está respaldado solo en parte:
  la fuerza estacional es 0.51 y el índice de 2024 y 2025 es casi
  idéntico (Figura 2), pero la ACF en el rezago 12 es 0.16.
\item
  \emph{Estacionariedad tras diferenciar.} ADF y KPSS coinciden en que
  \((1-B^{12})y_t\) es estacionaria (sección 4).
\item
  \emph{Innovaciones no correlacionadas y aproximadamente normales.} Se
  verifica con Ljung-Box. Es razonable porque los conteos mensuales son
  de miles.
\item
  \emph{Calendario sin ajustar.} Los meses tienen entre 20 y 23 días
  hábiles (alrededor de ±7 \%); el modelo principal no lo corrige y se
  prueba como sensibilidad.
\item
  \emph{Sin cambio estructural.} El lanzamiento y la pausa del acero no
  se modelan; su efecto se evalúa aparte.
\end{itemize}

\textbf{Métricas de éxito.} Con \(e_t = y_t - \hat y^{+}_t\) y \(n\)
pronósticos,

\[\text{MAE} = \frac{1}{n}\sum |e_t|, \qquad \text{RMSE} = \sqrt{\frac{1}{n}\sum e_t^2}, \qquad \text{MASE} = \frac{\text{MAE}}{\frac{1}{n_0 - 12}\sum_{t=13}^{n_0} |y_t - y_{t-12}|},\]

donde el denominador es el error del ingenuo estacional dentro del
entrenamiento (\(n_0\) meses). Un MASE menor que 1 indica que el modelo
supera al ingenuo del entrenamiento (Hyndman y Koehler, 2006). Se
calcula por ventana y se promedia.

\textbf{Validación.} Ventana expansiva sin barajar: el primer
entrenamiento va de enero de 2023 a diciembre de 2024 (24 meses) y
pronostica tres meses; se agrega un mes por ventana hasta pronosticar
julio a septiembre de 2026. Resultan 19 ventanas y 57 pronósticos. Las
referencias son el \textbf{ingenuo estacional} (SNAIVE: repite el mismo
mes del año anterior) y \textbf{AutoETS}. \textbf{Criterio de éxito:}
MASE menor que 1 y menor que el de SNAIVE en la mayoría de las ventanas,
confirmado con la prueba de Wilcoxon sobre los MASE pareados. Como el
lanzamiento del acero afecta las ventanas recientes, los resultados se
reportan por separado según el origen: antes de marzo de 2026 (15
ventanas, pronostican hasta mayo) y desde marzo de 2026 (4 ventanas,
pronostican de abril a septiembre).

\section{6. Primer modelo aplicado}\label{primer-modelo-aplicado}

El SARIMA se ajustó y evaluó en el notebook (Python 3,
\texttt{statsforecast} 2.1.1, \texttt{statsmodels} 0.14.6, semilla
2026). A lo largo de las 19 ventanas, AutoARIMA eligió
SARIMA\((1,1,0)(0,1,0)_{12}\) en 10, \((3,1,0)(0,1,0)_{12}\) en 7 y
otros dos órdenes en las 2 restantes.

{\def\LTcaptype{none} % do not increment counter
\begin{longtable}[]{@{}
  >{\raggedright\arraybackslash}p{(\linewidth - 10\tabcolsep) * \real{0.1667}}
  >{\raggedright\arraybackslash}p{(\linewidth - 10\tabcolsep) * \real{0.1667}}
  >{\raggedright\arraybackslash}p{(\linewidth - 10\tabcolsep) * \real{0.1667}}
  >{\raggedright\arraybackslash}p{(\linewidth - 10\tabcolsep) * \real{0.1667}}
  >{\raggedright\arraybackslash}p{(\linewidth - 10\tabcolsep) * \real{0.1667}}
  >{\raggedright\arraybackslash}p{(\linewidth - 10\tabcolsep) * \real{0.1667}}@{}}
\toprule\noalign{}
\begin{minipage}[b]{\linewidth}\raggedright
Modelo
\end{minipage} & \begin{minipage}[b]{\linewidth}\raggedright
MAE
\end{minipage} & \begin{minipage}[b]{\linewidth}\raggedright
RMSE
\end{minipage} & \begin{minipage}[b]{\linewidth}\raggedright
MASE, todas (19)
\end{minipage} & \begin{minipage}[b]{\linewidth}\raggedright
MASE, antes de mar-2026 (15)
\end{minipage} & \begin{minipage}[b]{\linewidth}\raggedright
MASE, desde mar-2026 (4)
\end{minipage} \\
\midrule\noalign{}
\endhead
\bottomrule\noalign{}
\endlastfoot
Ingenuo estacional (SNAIVE) & 248.9 & 349.6 & \textbf{0.59} &
\textbf{0.33} & 1.57 \\
AutoETS & 343.1 & 486.0 & 0.83 & 0.49 & 2.11 \\
SARIMA(1,1,0)(0,1,0)₁₂ & 279.6 & 413.7 & 0.62 & 0.43 & \textbf{1.36} \\
\end{longtable}
}

Tabla 2. Desempeño de pronóstico a 3 meses. MAE y RMSE en órdenes por
mes, sobre los 57 pronósticos.

\begin{figure}
\centering
\includegraphics[width=0.8\linewidth,height=\textheight,keepaspectratio,alt={Figura 4. Arriba, MASE por ventana (línea punteada: marzo de 2026). Abajo, pronóstico emitido tres meses antes contra lo observado.}]{./media-564a31000ad1b7c6/figuras/fig4_validacion.png}
\caption{Figura 4. Arriba, MASE por ventana (línea punteada: marzo de
2026). Abajo, pronóstico emitido tres meses antes contra lo observado.}
\end{figure}

\textbf{Lectura de los resultados.}

\begin{itemize}
\tightlist
\item
  \emph{El SARIMA no supera al ingenuo estacional.} Su MASE promedio es
  0.62 contra 0.59; gana en 9 de las 19 ventanas y Wilcoxon no encuentra
  diferencia (p = 0.57). Cumple la primera parte del criterio (MASE
  \textless{} 1) pero no la segunda. AutoETS es peor (0.83) y solo
  supera a SNAIVE en 4 ventanas.
\item
  \emph{Antes del lanzamiento del acero, repetir el mes del año anterior
  es lo mejor} (0.33 contra 0.43), porque el patrón de 2024 y 2025 se
  repite casi igual.
\item
  \emph{Con origen desde marzo de 2026, los tres modelos fallan} (MASE
  entre 1.36 y 2.11); el SARIMA es el menos malo, pero son solo 4
  ventanas. SNAIVE subestima 9 \% (Figura 4, abajo: sigue el nivel del
  año anterior); el SARIMA, 2 \%, pero para septiembre de 2026
  pronosticó cerca de 5,900 órdenes (Figura 4, abajo) contra 4,095
  observadas.
\item
  \emph{El orden del modelo es inestable} (cuatro órdenes distintos),
  esperable con 24 a 42 meses y una diferencia estacional que consume
  12.
\item
  \emph{Los residuos son compatibles con ruido blanco} (Ljung-Box p =
  0.71 con 6 rezagos y 0.61 con 12; \(n = 32\)), aunque con tan pocos
  residuos la prueba detecta poco. Su media es −79 y su desviación 397
  órdenes.
\item
  \emph{Sensibilidad por días hábiles.} Modelar el promedio por día
  hábil y multiplicar de vuelta mejora a los tres modelos (MASE: SNAIVE
  0.51, SARIMA 0.56, AutoETS 0.64) y no cambia el mejor del conjunto;
  con origen desde marzo de 2026 los resultados son AutoETS 0.98, SARIMA
  1.04 y SNAIVE 1.51.
\end{itemize}

{\def\LTcaptype{none} % do not increment counter
\begin{longtable}[]{@{}llll@{}}
\toprule\noalign{}
Mes & SARIMA (intervalo 80 \%) & SNAIVE & AutoETS \\
\midrule\noalign{}
\endhead
\bottomrule\noalign{}
\endlastfoot
Octubre 2026 & 5,172 (4,653 a 5,691) & 4,948 & 4,846 \\
Noviembre 2026 & 3,976 (3,416 a 4,536) & 4,187 & 4,847 \\
Diciembre 2026 & 4,204 (3,519 a 4,888) & 4,156 & 4,847 \\
\textbf{Total} & \textbf{13,352} & \textbf{13,291} & \textbf{14,540} \\
\end{longtable}
}

Tabla 3. Pronóstico de octubre a diciembre de 2026 con los 45 meses.

El total del trimestre es casi el mismo con SARIMA y con el ingenuo
(13,352 y 13,291), pero \textbf{el valor de octubre no es confiable}: al
día 25 ya hay 3,036 órdenes, y si los seis días restantes se parecen a
la última semana el mes cerraría cerca de 4,000, por debajo del
intervalo del 80 \%. El SARIMA arrastra el nivel alto de junio a agosto.
En resumen, con 45 meses el modelo estadístico clásico apenas iguala a
la regla más simple; el valor del ejercicio es mostrar dónde falla y por
qué.

\section{7. Próximos pasos}\label{pruxf3ximos-pasos}

Para la entrega final se probará un segundo método que atienda lo que
falló: el cambio de nivel del acero y los pocos datos. Se propone un
\textbf{modelo de fundación (Chronos-2)}: se usa sin entrenamiento, lo
que conviene con 45 meses, y admite covariables (Ansari et al., 2024),
como la proporción de órdenes de acero y los días hábiles. Se evaluará
con las mismas 19 ventanas y el mismo MASE, y se compararán con Prophet
(Taylor y Letham, 2018) como contraste y con el ajuste por días hábiles
de la sección 6. A mediano plazo se pasará al \textbf{pronóstico
diario}: una exploración preliminar con 1,394 días (notebook aparte)
muestra una estacionalidad semanal fuerte (0.89), de lunes a jueves unas
200 a 230 órdenes por día y fines de semana cerca de 25, y caídas
marcadas en feriados de EE. UU., que el modelo mensual no puede
capturar. También se confirmará con operaciones cómo se reabre el
producto de acero, porque cambia el nivel esperado de los próximos
meses.

\section{Declaración sobre el uso de inteligencia
artificial}\label{declaraciuxf3n-sobre-el-uso-de-inteligencia-artificial}

Se utilizó Claude (Anthropic), mediante Claude Code, para consultar la
base de datos de la empresa y exportar la serie, escribir y depurar el
código del notebook, y redactar un borrador de este documento a partir
de las instrucciones del avance y de los resúmenes de las clases 5 a 7.
Los resultados numéricos provienen de la ejecución del notebook adjunto.
\emph{(Completar: el grupo revisó, verificó y puede explicar cada parte
del documento y del código; indicar aquí quién hizo qué.)}

\section{Referencias}\label{referencias}

Ansari, A. F., Stella, L., Turkmen, C., Zhang, X., Mercado, P., Shen,
H., Shchur, O., Rangapuram, S. S., Arango, S. P., Kapoor, S.,
Zschiegner, J., Maddix, D. C., Wang, H., Mahoney, M. W., Torkkola, K.,
Wilson, A. G., Bohlke-Schneider, M., \& Wang, Y. (2024). Chronos:
Learning the language of time series. \emph{Transactions on Machine
Learning Research}. https://arxiv.org/abs/2403.07815

Cleveland, R. B., Cleveland, W. S., McRae, J. E., \& Terpenning, I.
(1990). STL: A seasonal-trend decomposition procedure based on loess.
\emph{Journal of Official Statistics, 6}(1), 3–73.

Dickey, D. A., \& Fuller, W. A. (1979). Distribution of the estimators
for autoregressive time series with a unit root. \emph{Journal of the
American Statistical Association, 74}(366), 427–431.
https://doi.org/10.1080/01621459.1979.10482531

Garza, A., Mergenthaler-Canseco, M., Challú, C., \& Olivares, K. G.
(2022). \emph{StatsForecast: Lightning fast forecasting with statistical
and econometric models}. PyCon Salt Lake City.
https://github.com/Nixtla/statsforecast

Hyndman, R. J., \& Athanasopoulos, G. (2021). \emph{Forecasting:
Principles and practice} (3.ª ed.). OTexts. https://otexts.com/fpp3/

Hyndman, R. J., \& Khandakar, Y. (2008). Automatic time series
forecasting: The forecast package for R. \emph{Journal of Statistical
Software, 27}(3), 1–22. https://doi.org/10.18637/jss.v027.i03

Hyndman, R. J., \& Koehler, A. B. (2006). Another look at measures of
forecast accuracy. \emph{International Journal of Forecasting, 22}(4),
679–688. https://doi.org/10.1016/j.ijforecast.2006.03.001

Kwiatkowski, D., Phillips, P. C. B., Schmidt, P., \& Shin, Y. (1992).
Testing the null hypothesis of stationarity against the alternative of a
unit root. \emph{Journal of Econometrics, 54}(1–3), 159–178.
https://doi.org/10.1016/0304-4076(92)90104-Y

LightForce Orthodontics. (2026). \emph{Tabla
analytics.manufacturing.mfg\_all\_orders} {[}Base de datos interna de
manufactura, Databricks{]}. Datos consultados el 26 de octubre de 2026
con autorización de la empresa.

Ljung, G. M., \& Box, G. E. P. (1978). On a measure of lack of fit in
time series models. \emph{Biometrika, 65}(2), 297–303.
https://doi.org/10.1093/biomet/65.2.297

Seabold, S., \& Perktold, J. (2010). statsmodels: Econometric and
statistical modeling with Python. \emph{Proceedings of the 9th Python in
Science Conference}, 92–96. https://doi.org/10.25080/Majora-92bf1922-011

Taylor, S. J., \& Letham, B. (2018). Forecasting at scale. \emph{The
American Statistician, 72}(1), 37–45.
https://doi.org/10.1080/00031305.2017.1380080

Wilcoxon, F. (1945). Individual comparisons by ranking methods.
\emph{Biometrics Bulletin, 1}(6), 80–83. https://doi.org/10.2307/3001968

\end{document}
